\section{The Conformal Subgraph Problem}
\label{sec:statement}
Both the distribution-free (Sections~\ref{sec:free} and~\ref{sec:fixed}) and distribution-based (Section~\ref{sec:distr2}) formulations lead to the same underlying optimization problem: given a weighted family of candidate routes with total weight~$W$, find a subgraph $G'\subseteq V$ that retains at least a $\tau = 1-\epsilon$ fraction of~$W$ while minimizing the number of edges (or another cost). The only difference lies in how the distribution over routes is obtained, via filtering on distance and assigning a utility, as in the distribution-free case, or directly from the model probabilities $g(\cdot\mid A)$ in the distribution-based case. 

The subgraph we output aims to approximate the smallest subgraph achieving the desired conformal guarantee in the perfectly calibrated case. We additionally require the algorithm to satisfy monotonicity (Definition~\ref{def:monotone}) to allow for valid calibration in the distribution-free setting.

\paragraph{Problem Statement.} The  conformal compression framework can therefore be expressed in combinatorial terms as follows. Each distinct candidate route $B$ in the historical data corresponds to a \emph{hyperedge} whose weight encodes its score or model probability.  The vertex set $V$ represents all atomic road segments (road edges or intersections) that may appear in any route, and each hyperedge $e\subseteq V$ corresponds to the subset of segments traversed by a particular route~$B$. Let $\mathcal E$ denote the set of hyperedges. The nonnegative weight $w_e$ of hyperedge~$e$ equals its score or its model probability, and let the total weight be $W=\sum_{e\in\mathcal E}w_e$.

Selecting a conformal subgraph that captures at least a fraction $\tau$ of this mass is therefore equivalent to selecting a subset of vertices $K\subseteq V$ that covers  hyper-edge weight $e(K)\ge \tau \cdot W$, where $e(S) := \sum_{\{e\in\mathcal E:\,e\subseteq S\}} w_e$. The parameter $\epsilon=1-\tau$ measures the fraction of probability mass allowed to remain uncovered, and the size of the chosen hypergraph vertex set reflects the number of edges in the compressed subgraph.

This abstraction leads to the following general formulation of the \emph{conformal subgraph problem}. Let $H=(V,\mathcal E)$ be a (possibly nonuniform) hypergraph with $|V|=n$ vertices and $|\mathcal E|=m$ hyperedges. Each hyperedge $e$ has a nonnegative weight $w_e$, and $W = \sum_{e\in\mathcal E}w_e$. Assume there exists an optimal vertex set $O\subseteq V$ satisfying
\[
  |O| = r,
  \qquad
  e(O) \ge W - q,
  \qquad
  q = \epsilon W,
\]
so that $O$ covers all but an $\epsilon$ fraction of the total weight.

We show in Appendix~\ref{app:np-hard} that the above problem is {\sc NP-Hard}, even in the regime where $\epsilon$ is a constant.  The goal is to approximate $O$ as closely as possible while allowing bicriteria trade-offs in coverage and size.

\begin{definition}[Bicriteria approximation]
\label{def:bicrit}
Given parameters $(\alpha,\beta)$, a vertex set $K\subseteq V$ is an \emph{$(\alpha,\beta)$-bicriteria approximation} if $
  W - e(K) \le \alpha\cdot \epsilon W$ and   $|K| \le \beta\,|O| = \beta \cdot r.$
\end{definition}

%In this formulation, the conformal calibration step determines the hypergraph structure and weights, and the conformal compression step corresponds to finding a vertex subset $K$ that approximately minimizes size while retaining almost all of the weighted probability mass. This leads to an approximately optimal compression had the weights corresponded to perfect calibration.

\noindent\textbf{Monotonicity.} While Definition~\ref{def:bicrit} formulates the problem for a fixed target mass $\tau = 1-\epsilon$, the distribution-free conformal framework (Section~\ref{sec:model}) requires a nested sequence of subgraphs (see Definition~\ref{def:monotone}). Therefore, we need our algorithm to additionally produce a family of solutions $\{K_\tau\}$ such that $K_{\tau_1} \subseteq K_{\tau_2}$ whenever $\tau_1 < \tau_2$. As shown in the example below, such a statement need not be true for the optimal solution, and one of our contributions is to show an approximation algorithm that is monotone.

\begin{example}
\label{eg:monotone}
Consider a setting with $3$ disjoint paths $P_1, P_2, P_3$ between $s$ and $t$. Paths $P_1$ and $P_2$ have two edges each, while path $P_3$ has three edges. Suppose $30\%$ of the traffic uses paths $P_1$, $30\%$ uses path $P_2$, and $40\%$ uses path $P_3$. If the desired coverage is $0.6$, the optimal solution chooses paths $P_1$ and $P_2$, while if it is $0.7$, the chosen paths are $P_1$ and $P_3$, so that the optimal solution is not monotone. 
\end{example}


